\subsection{CENTRALIZERS}

Recall that two elements \(x, y\) of a group are called permutable if \((x, y) = e\) , or \((\operatorname{Int} x)y = y\) , or \((\operatorname{Int} y)x = x\) ; and that two elements \(a, b\) of a Lie algebra are called permutable if \([a, b] = 0\) , or \((\operatorname{ad} a).b = 0\) , or \((\operatorname{ad} b).a = 0\) . Let \(G\) be a Lie group, \(x \in G\) , \(a \in L(G)\) ; \(x\) and \(a\) are called permutable if \((\operatorname{Ad} x).a = a\) , that is if \(xa = ax\) in \(T(G)\) .

Let \(G\) be a Lie group, \(g\) its Lie algebra, A a subset of \(G\) and \(a\) a subset of \(g\) . Let \(Z_{G}(A)\) (resp. \(Z_{G}(a)\) ) denote the set of elements of \(G\) which are permutable with all the elements of A (resp. \(a\) ). It is a closed subgroup of \(G\) . Let \(\mathfrak{z}_{\mathfrak{g}}(A)\) (resp. \(\mathfrak{z}_{\mathfrak{g}}(a)\) ) denote the set of elements of \(g\) which are permutable with all the elements of A (resp. \(a\) ). It is a closed Lie subalgebra of \(g\) .

PROPOSITION 7. Let \(G\) be a finite-dimensional Lie group, \(g\) its Lie algebra and \(a\) a subset of \(g\) . Then \(Z_G(a)\) is a Lie subgroup of \(G\) with Lie algebra \(\mathfrak{z}_{\mathfrak{g}}(a)\) .

This follows from § 3, Proposition 44 and Corollary 2 to Proposition 39.

PROPOSITION 8. Let \(\mathbf{G}\) be a finite-dimensional real or complex Lie group, \(\mathfrak{g}\) its Lie algebra and \(\mathbf{A}\) a subset of \(\mathbf{G}\) . Then \(Z_{\mathbf{G}}(\mathbf{A})\) is a Lie subgroup of \(\mathbf{G}\) with Lie algebra \(\mathfrak{z}_{\mathfrak{g}}(\mathbf{A})\) .

Suppose that A consists of a single point \(a\) . Then \(Z_{\mathbf{G}}(\mathbf{A})\) is the set of fixed points of Int \(a\) ; hence \(Z_{\mathbf{G}}(\mathbf{A})\) is a Lie subgroup of \(\mathbf{G}\) and \(\mathbf{L}(Z_{\mathbf{G}}(\mathbf{A}))\) is the set of fixed points of Ad \(a\) , that is \(\mathfrak{z}_{\mathfrak{g}}(\mathbf{A})\) (§ 3, no. 8, Corollary 1 to Proposition 29). The general case follows by means of § 6, no. 2, Corollary 3 to Proposition 1.

PROPOSITION 9. Let G be a finite-dimensional real or complex Lie group, g its Lie algebra, A an integral subgroup of G and \(\mathfrak{a} = \mathrm{L}(\mathbf{A})\) . Then \(\mathrm{Z}_{\mathbf{G}}(\mathbf{A}) = \mathrm{Z}_{\mathbf{G}}(\mathfrak{a})\) , \(\mathfrak{z}_{\mathfrak{g}}(\mathbf{A}) = \mathfrak{z}_{\mathfrak{g}}(\mathfrak{a})\) and \(\mathrm{Z}_{\mathbf{G}}(\mathbf{A})\) is a Lie subgroup of G with Lie algebra \(\mathfrak{z}_{\mathfrak{g}}(\mathfrak{a})\) . Let \(x \in G\) . Then

\[
\begin{array}{l l} x \in Z _ {\mathrm{G}} (\mathrm{A}) \Leftrightarrow \mathrm{A} \subset Z _ {\mathrm{G}} (\{x \}) \\ \quad \Leftrightarrow \mathfrak {a} \subset \mathrm{L} (Z _ {\mathrm{G}} (\{x \})) & (\S 6, \text { Corollary   2   to   Proposition   3 }) \\ \quad \Leftrightarrow \mathfrak {a} \subset \mathfrak {z}_ {\mathfrak {g}} (\{x \}) & (\text { Proposition   8 }) \\ \quad \Leftrightarrow x \in Z _ {\mathrm{G}} (\mathfrak {a}) \end{array}
\]

and hence \(Z_{\mathbf{G}}(\mathbf{A}) = Z_{\mathbf{G}}(\mathfrak{a})\) . Let \(u \in \mathfrak{g}\) . Then

\[
\begin{array}{r l} u \in \mathfrak{z}_{\mathfrak{g}}(\mathfrak{a}) & \Leftrightarrow \mathfrak{a} \subset Z_{\mathrm{G}} (\{u\}) \\ & \Leftrightarrow \mathfrak{a} \subset \mathrm{L} (Z_{\mathrm{G}} (\{u\})) \quad (\S 6, \text { Corollary   2   to   Proposition   3 }) \\ & \Leftrightarrow \mathfrak{a} \subset \mathfrak{z}_{\mathfrak{g}} (\{u\}) \quad (\text { Proposition   7 }) \\ & \Leftrightarrow u \in \mathfrak{z}_{\mathfrak{g}} (\mathfrak{a}) \end{array}
\]

and hence \(\mathfrak{z}_{\mathfrak{g}}(\mathbf{A}) = \mathfrak{z}_{\mathfrak{g}}(\mathfrak{a})\) . The last assertion then follows from Proposition 7 or Proposition 8.

